\documentclass[aspectratio=169]{beamer}
\input{header}
\coursecode{JKSSB}

\title{Office Shortcut \& Product Trap Lab}
\subtitle{Lesson 54 --- Exact Keys and Product-Category Traps}
\author{JKSSB Computer Awareness}
\date{\today}

\begin{document}
\frame{\titlepage}

\begin{frame}{This Lesson Is a Trap Drill}
  \begin{itemize}
    \item The JKSSB computer section tests \key{exact keys} and
      \key{exact product categories}, not vague recognition.
    \item A wrong option is usually the \key{nearest neighbour} --- the key
      or product that ``sounds right''.
    \item In this lesson every question is built around a real near-neighbour
      or defective-option pattern.
    \item \key{Read each option as a claim to be tested}, not merely as a
      phrase to recognise.
    \item \muted{Preserve the exact English term: F5, F7, Ctrl+H, \$A\$1.}
  \end{itemize}
  \begin{alertblock}{Trap drill}
    Same key, different program, different job --- and the same product name
    placed in the wrong software category.
  \end{alertblock}
\end{frame}

\begin{frame}{MS Word: Exact Shortcut Keys}
  \begin{center}
    \begin{tabular}{p{0.24\textwidth}p{0.62\textwidth}}
      \toprule
      \textbf{Shortcut} & \textbf{Action in MS Word} \\
      \midrule
      F7 & Spelling \& Grammar (spell check) \\
      Shift+F7 & Thesaurus \\
      Ctrl+H & Find and Replace \\
      Ctrl+F & Find (Navigation) \\
      Ctrl+K & Insert Hyperlink \\
      Ctrl+X / Ctrl+C / Ctrl+V & Cut / Copy / Paste \\
      \bottomrule
    \end{tabular}
  \end{center}
  \vspace{0.25cm}
  \muted{Ctrl+H always replaces; Ctrl+F only finds.}
\end{frame}

\begin{frame}{Word Views and Review Tools}
  \begin{itemize}
    \item \key{Print Layout} shows the document exactly as it will appear on
      the printed page.
    \item Other views: Web Layout, Outline, Draft.
    \item \key{Track Changes} records every insertion, deletion, and
      formatting edit.
    \item Each change can be \key{accepted or rejected} later.
    \item \muted{Track Changes also stores author and time as document
      metadata.}
  \end{itemize}
  \begin{alertblock}{Trap alert}
    Track Changes does not delete the original text permanently. The original
    is preserved until a reviewer accepts or rejects the change (PYQ-12).
  \end{alertblock}
\end{frame}

\begin{frame}{Word Trap Alert: Near-Neighbour Keys}
  \begin{itemize}
    \item \key{Ctrl+H} opens Find and Replace; \key{Ctrl+F} opens Find only.
    \item \key{F7} is spell check; \key{Shift+F7} is the Thesaurus.
    \item \key{F1} is Help; \key{F2} renames (in File Explorer) or edits a
      cell (in Excel).
    \item \key{Ctrl+K} inserts a hyperlink, not a comment.
    \item \muted{One letter or one modifier changes the whole action.}
  \end{itemize}
  \begin{block}{Habit to build}
    Say the full action with the key: ``F7 = Spelling \& Grammar; Ctrl+H =
    Find and Replace.''
  \end{block}
\end{frame}

\begin{frame}{Word: Section Break vs Page Break}
  \begin{itemize}
    \item A \key{Page Break} begins a new page but keeps the same section.
    \item A \key{Section Break} begins a new \key{section}.
    \item Only a \key{Section Break} allows a \key{different header or
      footer} on the following pages.
    \item \muted{Common wrong option: ``Page Break changes the
      header/footer.'' It does not (PYQ-13, TRAP-13).}
  \end{itemize}
  \begin{alertblock}{Trap alert}
    Different headers/footers require a \key{Section Break}, never a Page
    Break.
  \end{alertblock}
\end{frame}

\begin{frame}{PowerPoint: The Show Keys}
  \begin{center}
    \begin{tabular}{p{0.26\textwidth}p{0.60\textwidth}}
      \toprule
      \textbf{Shortcut} & \textbf{Action in PowerPoint} \\
      \midrule
      F5 & Start the slide show from the \textbf{beginning} \\
      Shift+F5 & Start the slide show from the \textbf{current} slide \\
      Alt+F5 & Start the presentation in \textbf{Presenter View} \\
      Esc & \textbf{End} the running slide show \\
      \bottomrule
    \end{tabular}
  \end{center}
  \vspace{0.25cm}
  \muted{F5 begins; Shift+F5 resumes from here; Alt+F5 is Presenter View;
  Esc ends (PYQ-40).}
\end{frame}

\begin{frame}{PowerPoint Trap Alert: F5 Is Not Spell Check}
  \begin{itemize}
    \item In PowerPoint, \key{F5} starts the slide show.
    \item In Word, \key{F7} is spell check; in PowerPoint too, F7 checks
      spelling.
    \item \key{F5} does \key{not} spell-check, and \key{F7} does \key{not}
      start a show.
    \item \key{Esc} ends the show; it does not go to the first slide.
    \item \muted{The same key can mean different things in different
      programs.}
  \end{itemize}
  \begin{alertblock}{Trap alert}
    ``F5 = spell check'' is the classic near-neighbour trap. F5 starts the
    show; spell check is F7 (TRAP-32).
  \end{alertblock}
\end{frame}

\begin{frame}{Handouts, Slide Master, Orientation}
  \begin{itemize}
    \item \key{Handouts} print several slides per page for the audience, with
      optional note lines.
    \item \key{Handout Master} controls the handout and notes-page layout.
    \item The \key{Slide Master} controls the design, fonts, colours, and
      layout of \key{all slides} that use it.
    \item A change to the Slide Master propagates to every dependent slide.
    \item \key{Orientation} switches between \key{Portrait} and
      \key{Landscape}, for slides as well as for notes and handouts.
  \end{itemize}
\end{frame}

\begin{frame}{PowerPoint: SmartArt Conversion}
  \begin{itemize}
    \item \key{SmartArt} turns text into a visual diagram (process, cycle,
      hierarchy, and more).
    \item A \key{bulleted list}, a \key{block of text}, and \key{indented
      text} can all be converted to SmartArt.
    \item A \key{table of data} \key{cannot} be directly converted to
      SmartArt.
    \item To move table data into SmartArt, the text must be copied into a
      text box first.
    \item \muted{Content type decides whether conversion is allowed
      (PYQ-17).}
  \end{itemize}
  \begin{alertblock}{Trap alert}
    ``Convert to SmartArt'' is disabled for a table --- a table is not a
    simple text string.
  \end{alertblock}
\end{frame}

\begin{frame}{Excel: The Formula Palette}
  \begin{itemize}
    \item A formula begins with \key{=}; without it, Excel treats the entry
      as text (PYQ-38, TRAP-30).
    \item The \key{formula palette} (the Insert Function dialog) helps you
      \key{build a formula}.
    \item It lists the function and shows each \key{argument} in its own
      box.
    \item It can be opened from the \key{fx} button beside the formula bar.
    \item \muted{It guides formula entry; it does not change cell size or
      data.}
  \end{itemize}
  \begin{block}{Definition}
    The formula palette is the dialog that lets you pick a function and fill
    its arguments while building a formula.
  \end{block}
\end{frame}

\begin{frame}{Excel: Cell References and the \$ Sign}
  \begin{center}
    \begin{tabular}{p{0.20\textwidth}p{0.32\textwidth}p{0.36\textwidth}}
      \toprule
      \textbf{Reference} & \textbf{Type} & \textbf{Behaviour when copied} \\
      \midrule
      A1 & Relative & Both column and row change \\
      \$A\$1 & Absolute & Neither column nor row changes \\
      \$A1 & Mixed & Column fixed, row changes \\
      A\$1 & Mixed & Row fixed, column changes \\
      \bottomrule
    \end{tabular}
  \end{center}
  \vspace{0.25cm}
  \muted{The \$ sign ``locks'' the part that follows it. Press F4 while
  editing to cycle the reference type.}
\end{frame}

\begin{frame}{Excel Trap Alert: Locking and Starting a Formula}
  \begin{itemize}
    \item \key{\$A\$1} is an \key{absolute} reference: the column and the row
      are both locked.
    \item \key{A1} is relative; it shifts when the formula is copied.
    \item A formula must start with \key{=}; \key{+} or \key{\%} is not the
      formula marker (PYQ-38).
    \item \key{F4} toggles relative, absolute, and mixed references.
    \item \muted{\$A1 and A\$1 are mixed --- only one part is locked.}
  \end{itemize}
  \begin{alertblock}{Trap alert}
    A single \$ both locks a reference and can be forgotten --- \key{\$A\$1}
    means the reference never moves when copied.
  \end{alertblock}
\end{frame}

\begin{frame}{Product-Category Traps: Word, Excel, Access, MS-DOS}
  \begin{center}
    \begin{tabular}{p{0.26\textwidth}p{0.58\textwidth}}
      \toprule
      \textbf{Program} & \textbf{Correct category} \\
      \midrule
      MS Word & Word processor (text documents) \\
      MS Excel & Spreadsheet (cells, formulas) \\
      MS Access & Database management system (tables, queries) \\
      MS-DOS & Disk Operating System (command-line OS) \\
      \bottomrule
    \end{tabular}
  \end{center}
  \vspace{0.25cm}
  \muted{MS Access is a database, not a spreadsheet. None of Word, Excel, or
  Access is an operating system.}
\end{frame}

\begin{frame}{Word Feature Trap: Mail Merge}
  \begin{itemize}
    \item \key{Mail Merge} combines a \key{main document} with a
      \key{data source} (such as an address database).
    \item It is a \key{Word} feature, used to produce many personalised
      letters, labels, or envelopes.
    \item It is \key{not} Track Changes, AutoCorrect, or a spell-check tool.
    \item \muted{Main document + data source = merged output (PYQ-37,
      TRAP-29).}
  \end{itemize}
  \begin{alertblock}{Trap alert}
    ``Combines letters with an address database'' = \key{Mail Merge}. Do not
    confuse it with Track Changes.
  \end{alertblock}
\end{frame}

\begin{frame}{Exact Shortcut Near-Neighbours: F1, F2, F5, F7}
  \begin{center}
    \begin{tabular}{p{0.12\textwidth}p{0.28\textwidth}p{0.46\textwidth}}
      \toprule
      \textbf{Key} & \textbf{In Word} & \textbf{In PowerPoint / Windows} \\
      \midrule
      F1 & Help & Help \\
      F2 & Move text / edit field & Rename (File Explorer); edit cell (Excel) \\
      F5 & Go To & Start slide show (PowerPoint) \\
      F7 & Spelling \& Grammar & Spelling (PowerPoint) \\
      \bottomrule
    \end{tabular}
  \end{center}
  \vspace{0.25cm}
  \muted{Ask ``which program?'' before you answer a function-key question.}
\end{frame}

\begin{frame}{Trap Alert: Terms to Keep Straight}
  \begin{itemize}
    \item Word: \key{F7} = spell check; \key{Ctrl+H} = Find and Replace;
      \key{Ctrl+K} = hyperlink.
    \item Word: \key{Section Break} (not Page Break) gives different
      headers/footers (TRAP-13).
    \item PowerPoint: \key{F5} begins, \key{Shift+F5} from current,
      \key{Alt+F5} Presenter View, \key{Esc} ends (TRAP-32).
    \item PowerPoint: a \key{table} cannot convert directly to SmartArt.
    \item Excel: \key{\$A\$1} is absolute; a formula must start with \key{=}.
    \item Products: Word = word processor, Excel = spreadsheet,
      Access = database, MS-DOS = operating system.
  \end{itemize}
\end{frame}

\begin{frame}{Lesson 54: Recall}
  \begin{itemize}
    \item Word keys: F7 spell check, Ctrl+H replace, Ctrl+K hyperlink,
      Ctrl+X/V cut and paste, Print Layout view, Track Changes for revisions.
    \item Section Break, not Page Break, enables different headers/footers.
    \item PowerPoint keys: F5 from the beginning, Shift+F5 from current,
      Alt+F5 Presenter View, Esc to end.
    \item Handouts print several slides per page; the Slide Master controls
      all slides; orientation is Portrait or Landscape.
    \item A table cannot be converted directly to SmartArt.
    \item Excel: the formula palette builds formulas; \key{\$A\$1} is an
      absolute reference.
    \item Product categories: Word processor vs spreadsheet vs database vs
      operating system.
  \end{itemize}
\end{frame}
\end{document}
